\section{Introduction}
\label{sec:intro}

% Story arc (draft prose; tighten after results are final).
A project assistant sits in on sixteen weekly meetings.
In week five, one engineer remarks that the nightly backup does not cover the audit-log database; nobody mentions it again.
Meanwhile the migration timeline is restated every week.
By week sixteen the assistant's notes describe the timeline in detail, and say nothing about the backup.
Nothing in the notes is false; they are fluent, on-topic and within budget.
What they have lost is the one claim that should change a go-live decision.

Long-running agents increasingly rely on exactly this kind of memory: a bounded text that a model rewrites after every interaction \citep{wang2023recursively,packer2023memgpt,chhikara2025mem0,xu2025amem}.
That repeated rewriting loses information is by now well documented, under the names context collapse \citep{zhang2026ace}, the compaction cliff \citep{zerhoudi2026compaction} and goal drift \citep{zahn2026facts}.
What is not known is \emph{which} information is lost and \emph{why}.
Observational evidence suggests that frequency matters: concerns raised once are the ones single-pass summaries drop \citep{tamme2027voice}, and summarisers behave like a majority vote \citep{tian2026rgmem}.
But in natural data a claim's frequency is confounded with its content, its speaker, its recency and its position, so these observations cannot say whether repetition itself protects a claim.

We isolate it.
\sys{} generates multi-session meeting dialogues from fixed claim banks in which each claim's content and paraphrases are fixed and counterbalanced across dialogues, and only its mention frequency $f\in\{1,2,4\}$ varies; the session of its last mention is drawn independently of $f$.
Because the same claim appears at every frequency in different dialogues, and because a full-context reader is probed on the same questions, we can estimate the effect of frequency that is specific to iterative summarisation as a difference-in-differences.

\paragraph{Contributions.}
\begin{itemize}
\item A controlled benchmark for claim survival in iterative memory, with four domains, decision-critical claim types (warnings, preconditions, dissent) and incidental controls, plus full-context, truncation and empty-memory baselines (\S\ref{sec:benchmark}).
\item Evidence that rolling summarisation erases single-mention critical claims at a much higher rate than repeated ones, while full context and truncation are flat in frequency, across \todo{N} model--domain cells; the effect survives a clustered schedule that removes the first-mention confound (\S\ref{sec:findings}).
\item \ledger{}, an extract-then-consolidate memory with tiered eviction, compared at equal budget against prompt-level, typed-compaction \citep{zerhoudi2026compaction} and delta-update \citep{zhang2026ace} baselines (\S\ref{sec:method}, \S\ref{sec:method-results}).
\end{itemize}

\todo{One paragraph of headline numbers once all runs are in: retention at f=1 vs f=4; ledger gain at f=1; when the prompt fix suffices (7B, loose budget).}
